% =============================================================================
% Introduction
% =============================================================================
\section{Introduction}\label{sec:introduction}

In 1966, Kac posed the question ``Can one hear the shape of a drum?''
\cite{Kac1966}, asking whether the eigenfrequencies of the Laplacian on a
planar domain uniquely determine its geometry.  Gordon, Webb, and Wolpert
\cite{Gordon1992} famously answered in the negative by constructing pairs of
non-isometric planar domains with identical Dirichlet spectra.  Yet the
question has continued to inspire inverse spectral problems across
mathematics, physics, and engineering.

In the present work we transpose Kac's question to the vibroacoustic setting:
given the measured resonant frequencies of a fluid-filled viscoelastic shell,
can one uniquely recover its geometric and material parameters?  This is not
merely an academic curiosity.  Non-invasive assessment of soft-tissue
mechanical properties underpins clinical elastography
\cite{Sarvazyan1998,Doyley2012}, and emerging vibroacoustic techniques seek
to characterise organs \emph{in vivo} from their spectral response.

We consider a class of models motivated by the human abdomen: a thin-walled,
pre-stressed, fluid-filled oblate spheroidal shell.  The forward problem ---
predicting modal frequencies from parameters --- is addressed in a companion
paper using a Rayleigh--Ritz variational formulation
\cite{BrowntoneCompanion}.  Here we study the \emph{inverse}: given $N$
measured flexural frequencies $f_2, f_3, \dots, f_{N+1}$, can we recover the
semi-major axis~$a$, semi-minor axis~$c$, and Young's modulus~$E$?

Our central finding is that identifiability depends critically on the
\emph{geometry} of the forward model:
\begin{itemize}
  \item For the \textbf{equivalent-sphere} approximation, the Jacobian of
    the frequency--parameter map is rank-deficient (condition number
    $\sim 10^{10}$), because all flexural modes scale identically with the
    effective radius.
  \item For the \textbf{oblate spheroid} (Ritz) model, the Jacobian has full
    rank (condition number $\approx 83.5$), enabling sub-percent parameter
    recovery under realistic noise.
\end{itemize}
One can, in a practical sense, ``hear the shape'' of an oblate organ --- but
not a spherical one.

We frame this as a problem of \emph{practical identifiability}
\cite{Raue2009,Wieland2021}, characterised by the condition number of the
Jacobian and the Cram\'er--Rao lower bounds on parameter uncertainty, rather
than claiming a uniqueness theorem that would require rigorous proof of global
injectivity.

The paper is organised as follows.  Section~\ref{sec:theory} summarises the
forward model.  Section~\ref{sec:methodology} develops the inverse
framework: Jacobian computation, Fisher information, and Cram\'er--Rao
bounds.  Section~\ref{sec:results} presents the sphere--oblate comparison,
parameter-space condition maps, and noise-sensitivity analysis.
Section~\ref{sec:discussion} discusses implications for tissue
characterisation, and Section~\ref{sec:conclusion} concludes.
